\begin{table}[H]
\centering
\caption{Transported Endline Effects Reweighted to the Baseline Frame}
\label{tab:transported_effects_t1}
\footnotesize
\begin{threeparttable}
\begin{tabular}{llcc}
\toprule
Outcome & Contrast & Internal ITT & Transported ITT \\
\midrule
Foreign media receptivity/demand index (endline) & Apolitical China & -0.004 & -0.004 \\
Foreign media receptivity/demand index (endline) & Pooled Anti-China & -0.091*** & -0.096*** \\
Foreign media receptivity/demand index (endline) & Pooled Pro-China & 0.223*** & 0.232*** \\
Identity backlash index (endline) & Apolitical China & 0.004 & 0.012 \\
Identity backlash index (endline) & Pooled Anti-China & 0.062*** & 0.074*** \\
Identity backlash index (endline) & Pooled Pro-China & -0.018 & -0.010 \\
Legitimacy index (endline) & Apolitical China & 0.042*** & 0.034** \\
Legitimacy index (endline) & Pooled Anti-China & -0.033** & -0.036** \\
Legitimacy index (endline) & Pooled Pro-China & 0.075*** & 0.072*** \\
\bottomrule
\end{tabular}
\begin{tablenotes}[flushleft]
\footnotesize
\item Entries compare the internal endline ITT estimates from the randomized study sample with transported estimates reweighted to the baseline respondent frame. Transport weights combine the inverse predicted probability of entering the randomized sample and the inverse predicted probability of completing the endline survey, both estimated from baseline observables; probabilities are trimmed to [0.05, 0.95] and the final weights are normalized to have mean one among endline respondents. Transporting beyond the randomized sample relies on selection-on-observables and is therefore descriptive rather than design-based.
\end{tablenotes}
\end{threeparttable}
\end{table}
